\documentclass[a4paper,12pt,twoside]{memoir}

% Castellano
\usepackage[spanish,es-tabla]{babel}
\selectlanguage{spanish}
\usepackage[utf8]{inputenc}
\usepackage[T1]{fontenc}
\usepackage{lmodern} % scalable font
\usepackage{microtype}
\usepackage{placeins}

\RequirePackage{booktabs}
\RequirePackage[table]{xcolor}
\RequirePackage{xtab}
\RequirePackage{multirow}

% Links
\usepackage[colorlinks]{hyperref}
\hypersetup{
	allcolors = {red}
}

% Ecuaciones
\usepackage{amsmath}

% Rutas de fichero / paquete
\newcommand{\ruta}[1]{{\sffamily #1}}

% Párrafos
\nonzeroparskip


% Imagenes
\usepackage{graphicx}
\newcommand{\imagen}[2]{
	\begin{figure}[!h]
		\centering
		\includegraphics[width=0.9\textwidth]{#1}
		\caption{#2}\label{fig:#1}
	\end{figure}
	\FloatBarrier
}

\newcommand{\imagenflotante}[2]{
	\begin{figure}%[!h]
		\centering
		\includegraphics[width=0.9\textwidth]{#1}
		\caption{#2}\label{fig:#1}
	\end{figure}
}



% El comando \figura nos permite insertar figuras comodamente, y utilizando
% siempre el mismo formato. Los parametros son:
% 1 -> Porcentaje del ancho de página que ocupará la figura (de 0 a 1)
% 2 --> Fichero de la imagen
% 3 --> Texto a pie de imagen
% 4 --> Etiqueta (label) para referencias
% 5 --> Opciones que queramos pasarle al \includegraphics
% 6 --> Opciones de posicionamiento a pasarle a \begin{figure}
\newcommand{\figuraConPosicion}[6]{%
  \setlength{\anchoFloat}{#1\textwidth}%
  \addtolength{\anchoFloat}{-4\fboxsep}%
  \setlength{\anchoFigura}{\anchoFloat}%
  \begin{figure}[#6]
    \begin{center}%
      \Ovalbox{%
        \begin{minipage}{\anchoFloat}%
          \begin{center}%
            \includegraphics[width=\anchoFigura,#5]{#2}%
            \caption{#3}%
            \label{#4}%
          \end{center}%
        \end{minipage}
      }%
    \end{center}%
  \end{figure}%
}

%
% Comando para incluir imágenes en formato apaisado (sin marco).
\newcommand{\figuraApaisadaSinMarco}[5]{%
  \begin{figure}%
    \begin{center}%
    \includegraphics[angle=90,height=#1\textheight,#5]{#2}%
    \caption{#3}%
    \label{#4}%
    \end{center}%
  \end{figure}%
}
% Para las tablas
\newcommand{\otoprule}{\midrule [\heavyrulewidth]}
%
% Nuevo comando para tablas pequeñas (menos de una página).
\newcommand{\tablaSmall}[5]{%
 \begin{table}
  \begin{center}
   \rowcolors {2}{gray!35}{}
   \begin{tabular}{#2}
    \toprule
    #4
    \otoprule
    #5
    \bottomrule
   \end{tabular}
   \caption{#1}
   \label{tabla:#3}
  \end{center}
 \end{table}
}

%
%Para el float H de tablaSmallSinColores
\usepackage{float}

%
% Nuevo comando para tablas pequeñas (menos de una página).
\newcommand{\tablaSmallSinColores}[5]{%
 \begin{table}[H]
  \begin{center}
   \begin{tabular}{#2}
    \toprule
    #4
    \otoprule
    #5
    \bottomrule
   \end{tabular}
   \caption{#1}
   \label{tabla:#3}
  \end{center}
 \end{table}
}

\newcommand{\tablaApaisadaSmall}[5]{%
\begin{landscape}
  \begin{table}
   \begin{center}
    \rowcolors {2}{gray!35}{}
    \begin{tabular}{#2}
     \toprule
     #4
     \otoprule
     #5
     \bottomrule
    \end{tabular}
    \caption{#1}
    \label{tabla:#3}
   \end{center}
  \end{table}
\end{landscape}
}

%
% Nuevo comando para tablas grandes con cabecera y filas alternas coloreadas en gris.
\newcommand{\tabla}[6]{%
  \begin{center}
    \tablefirsthead{
      \toprule
      #5
      \otoprule
    }
    \tablehead{
      \multicolumn{#3}{l}{\small\sl continúa desde la página anterior}\\
      \toprule
      #5
      \otoprule
    }
    \tabletail{
      \hline
      \multicolumn{#3}{r}{\small\sl continúa en la página siguiente}\\
    }
    \tablelasttail{
      \hline
    }
    \bottomcaption{#1}
    \rowcolors {2}{gray!35}{}
    \begin{xtabular}{#2}
      #6
      \bottomrule
    \end{xtabular}
    \label{tabla:#4}
  \end{center}
}

%
% Nuevo comando para tablas grandes con cabecera.
\newcommand{\tablaSinColores}[6]{%
  \begin{center}
    \tablefirsthead{
      \toprule
      #5
      \otoprule
    }
    \tablehead{
      \multicolumn{#3}{l}{\small\sl continúa desde la página anterior}\\
      \toprule
      #5
      \otoprule
    }
    \tabletail{
      \hline
      \multicolumn{#3}{r}{\small\sl continúa en la página siguiente}\\
    }
    \tablelasttail{
      \hline
    }
    \bottomcaption{#1}
    \begin{xtabular}{#2}
      #6
      \bottomrule
    \end{xtabular}
    \label{tabla:#4}
  \end{center}
}

%
% Nuevo comando para tablas grandes sin cabecera.
\newcommand{\tablaSinCabecera}[5]{%
  \begin{center}
    \tablefirsthead{
      \toprule
    }
    \tablehead{
      \multicolumn{#3}{l}{\small\sl continúa desde la página anterior}\\
      \hline
    }
    \tabletail{
      \hline
      \multicolumn{#3}{r}{\small\sl continúa en la página siguiente}\\
    }
    \tablelasttail{
      \hline
    }
    \bottomcaption{#1}
  \begin{xtabular}{#2}
    #5
   \bottomrule
  \end{xtabular}
  \label{tabla:#4}
  \end{center}
}



\definecolor{cgoLight}{HTML}{EEEEEE}
\definecolor{cgoExtralight}{HTML}{FFFFFF}

%
% Nuevo comando para tablas grandes sin cabecera.
\newcommand{\tablaSinCabeceraConBandas}[5]{%
  \begin{center}
    \tablefirsthead{
      \toprule
    }
    \tablehead{
      \multicolumn{#3}{l}{\small\sl continúa desde la página anterior}\\
      \hline
    }
    \tabletail{
      \hline
      \multicolumn{#3}{r}{\small\sl continúa en la página siguiente}\\
    }
    \tablelasttail{
      \hline
    }
    \bottomcaption{#1}
    \rowcolors[]{1}{cgoExtralight}{cgoLight}

  \begin{xtabular}{#2}
    #5
   \bottomrule
  \end{xtabular}
  \label{tabla:#4}
  \end{center}
}




\graphicspath{ {./img/} }

% Capítulos
\chapterstyle{bianchi}
\newcommand{\capitulo}[2]{
	\setcounter{chapter}{#1}
	\setcounter{section}{0}
	\chapter*{#2}
	\addcontentsline{toc}{chapter}{#2}
	\markboth{#2}{#2}
}

% Apéndices
\renewcommand{\appendixname}{Apéndice}
\renewcommand*\cftappendixname{\appendixname}

\newcommand{\apendice}[1]{
	%\renewcommand{\thechapter}{A}
	\chapter{#1}
}

\renewcommand*\cftappendixname{\appendixname\ }

% Formato de portada
\makeatletter
\usepackage{xcolor}
\newcommand{\tutor}[1]{\def\@tutor{#1}}
\newcommand{\course}[1]{\def\@course{#1}}
\definecolor{cpardoBox}{HTML}{E6E6FF}
\def\maketitle{
  \null
  \thispagestyle{empty}
  % Cabecera ----------------
\noindent\includegraphics[width=\textwidth]{cabecera}\vspace{1cm}%
  \vfill
  % Título proyecto y escudo informática ----------------
  \colorbox{cpardoBox}{%
    \begin{minipage}{.8\textwidth}
      \vspace{.5cm}\Large
      \begin{center}
      \textbf{TFG del Grado en Ingeniería Informática}\vspace{.6cm}\\
      \textbf{\LARGE\@title{}}
      \end{center}
      \vspace{.2cm}
    \end{minipage}

  }%
  \hfill\begin{minipage}{.20\textwidth}
    \includegraphics[width=\textwidth]{escudoInfor}
  \end{minipage}
  \vfill
  % Datos de alumno, curso y tutores ------------------
  \begin{center}%
  {%
    \noindent\LARGE
    Presentado por \@author{}\\ 
    en Universidad de Burgos --- \@date{}\\
    Tutor: \@tutor{}\\
  }%
  \end{center}%
  \null
  \cleardoublepage
  }
\makeatother


% Datos de portada
\title{título del TFG \\Documentación Técnica}
\author{nombre alumno}
\tutor{nombre tutor}
\date{\today}

\begin{document}

\maketitle



\cleardoublepage



%%%%%%%%%%%%%%%%%%%%%%%%%%%%%%%%%%%%%%%%%%%%%%%%%%%%%%%%%%%%%%%%%%%%%%%%%%%%%%%%%%%%%%%%



\frontmatter


\clearpage

% Indices
\tableofcontents

\clearpage

\listoffigures

\clearpage

\listoftables

\clearpage

\mainmatter

\appendix

\apendice{Plan de Proyecto Software}

\section{Introducción}

Una de las fases más importantes de cualquier proyecto es la planificación. En esta fase se estima el tiempo y los requisitos necesarios para la realización del proyecto. Para ello es necesaria una idea clara de lo que se quiere realizar de manera que todas las partes que componen el proyecto puedan ser analizadas individualmente.

Podemos dividir este anexo en dos apartados más pequeños:

\begin{itemize}
\item
  Planificación temporal.
\item
  Estudio de viabilidad.
\end{itemize}

En la primera fase se realiza una planificación de los tiempos esperados. Se especifican los tiempos necesarios para cada una de las partes que componen el proyecto, a la vez que una fecha de inicio y de fin de las mismas.

La segunda fase se centra en un estudio de la viabilidad del proyecto. En esta fase se pueden diferenciar dos partes:

\begin{itemize}
\item
\textbf{Viabilidad económica:} Estimación de los costes y beneficios del proyecto.
\item
\textbf{Viabilidad legal:} Regulaciones legales que pueden afectar al proyecto. En este caso serían la Ley de Protección de Datos y las licencias del software.
\end{itemize}

\section{Planificación temporal}

Para realizar una planificación del proyecto adecuada, se decidió aplicar \emph{Scrum} (una metodología ágil de desarrollo, explicada en la memoria).

Para ello, se dividió el trabajo en \emph{sprints} los cuales iban siendo planificados según se terminaba el anterior. La duración media de estos era de una semana de trabajo.

En cada uno de los \emph{sprint} se generan una serie de tareas las cuales tienen que ser realizadas en ese intervalo de tiempo. Para ello se utiliza la plataforma \emph{Zenhub}, en la cual las \emph{issues} hacen de tareas, y los \emph{Milestones} de \emph{Sprints}.

\subsection{Sprint 0}

Este primer \emph{sprint} fue el más complicado de planificar. Estaba incluido el desarrollo del \emph{hardware} el cual no sabíamos seguro si se conseguirían los conocimientos necesarios para diseñarlo, ni si después funcionaría correctamente. Además, el desarrollo de la aplicación dependía del resultado de esta parte ya que iba a basarse en ello.

Al final, sobrepasando un poco el tiempo planificado, se consiguió poner el \emph{hardware} en marcha, con lo que se podía comenzar con el desarrollo del \emph{software}.


\subsection{Sprint 1}

Los objetivos de este \emph{sprint} fueron asentar la idea general de la aplicación a desarrollar, su estructura principal y funcionalidades.

Además se comenzó a diseñar la estructura de la base de datos y el acceso al \emph{hardware} a través de la comunicación con los puertos serie.

Las horas de trabajo fueron sobrepasadas sobre lo planificado en el \emph{Sprint}.

\imagen{sprint1}{\emph{Sprint} 1.}

\subsection{Sprint 2}

Los objetivos de este \emph{sprint} fueron comenzar con el desarrollo de la interfaz gráfica. Para ello era necesario la creación de gráficas de forma dinámica. Estas serían utilizadas para monitorizar los datos que fluían por el bus CAN en tiempo real.

Esta fase llevó un poco más tiempo de lo esperado.

\imagen{sprint2}{\emph{Sprint} 2.}

\subsection{Sprint 3}

En este \emph{sprint} se planteó la posibilidad de poder importar y exportar un proyecto de la aplicación. Para ello, se decidió usar ficheros JSON, los cuales contendrían todos los datos del proyecto.

En este punto se planteó la posibilidad de realizar una parte de la aplicación sobre una tecnología web. La idea era generar los ficheros JSON desde la aplicación Java para posteriormente poder importarlos a la web para su monitorización desde la misma.

\imagen{sprint3}{\emph{Sprint} 3.}

\subsection{Sprint 4}

El objetivo de este \emph{sprint} era la generación de un gestor de proyectos interno en la aplicación. Para ello era necesario diseñar las funcionalidades para la creación de un nuevo proyecto, la asignación de valores identificados al mismo, así como sus opciones de guardar y eliminar de la base de datos.

\imagen{sprint4}{\emph{Sprint} 4.}

\subsection{Sprint 5}

En este \emph{sprint} se desarrolló la interacción con la base de datos, tanto para obtener datos de ella como para recuperarlos. Además se realizó un tratado de las posibles excepciones que puedan surgir por problemas internos de la aplicación.

\imagen{sprint5}{\emph{Sprint} 5.}

\subsection{Sprint 6}

En este \emph{sprint} se centraron los esfuerzos en tener un borrador de la memoria y los anexos. Además, el desarrollo de la aplicación continuaba añadiendo opciones como la de ver los datos de una ID y un \emph{byte} concretos, en tiempo real. Esta opción estaría disponible para todos los datos mostrados por pantalla.

\imagen{sprint6}{\emph{Sprint} 6.}

\subsection{Sprint 7}

En este último \emph{sprint} se trató de concluir con la documentación y terminar de pulir las pequeñas características que quedaban pendientes del proyecto. También se centraron esfuerzos en concluir la documentación y los vídeos de prueba de la aplicación.


\subsection{Resumen}

En la siguiente tabla se muestran los tiempos dedicados a cada uno de los distintos tipos de tareas:



\begin{longtable}[]{@{}lrr@{}}
\toprule
\begin{minipage}[b]{0.37\columnwidth}\raggedright\strut
Categoría\strut
\end{minipage} & \begin{minipage}[b]{0.37\columnwidth}\raggedleft\strut
Tiempo (horas)\strut
\end{minipage}\tabularnewline
\midrule
\endhead
\begin{minipage}[t]{0.37\columnwidth}\raggedright\strut
\emph{Documentación}\strut
\end{minipage} & \begin{minipage}[t]{0.37\columnwidth}\raggedleft\strut
67\strut
\end{minipage}\tabularnewline
\begin{minipage}[t]{0.37\columnwidth}\raggedright\strut
\emph{Características}\strut
\end{minipage} & \begin{minipage}[t]{0.37\columnwidth}\raggedleft\strut
132\strut
\end{minipage}\tabularnewline
\begin{minipage}[t]{0.37\columnwidth}\raggedright\strut
\emph{Investigación}\strut
\end{minipage}& \begin{minipage}[t]{0.37\columnwidth}\raggedleft\strut
42\strut
\end{minipage}\tabularnewline
\begin{minipage}[t]{0.37\columnwidth}\raggedright\strut
\emph{Corrección de errores}\strut
\end{minipage} & \begin{minipage}[t]{0.37\columnwidth}\raggedleft\strut
28\strut
\end{minipage}\tabularnewline
\midrule
\begin{minipage}[t]{0.37\columnwidth}\raggedright\strut
Total\strut
\end{minipage} & \begin{minipage}[t]{0.37\columnwidth}\raggedleft\strut
269\strut
\end{minipage}\tabularnewline
\bottomrule
\caption{Horas dedicadas al proyecto.}
\end{longtable}


\imagen{resumenHoras}{Resumen de horas dedicadas al proyecto.}


\section{Estudio de viabilidad}

En este punto se van a analizar por separado la viabilidad económica y la viabilidad legal del proyecto.

\subsection{Viabilidad económica}

\subsection{Costes}\label{costes}

Los costes van a ser desglosados en dos partes:

\textbf{Costes de personal:}

El proyecto ha sido desarrollado por una sola persona, empleada a tiempo completo durante tres meses. Considerando los siguientes valores:


\begin{longtable}[]{@{}lr@{}}
\toprule
\begin{minipage}[b]{0.38\columnwidth}\raggedright\strut
\textbf{Concepto}\strut
\end{minipage} & \begin{minipage}[b]{0.20\columnwidth}\raggedleft\strut
\textbf{Coste \euro{}}\strut
\end{minipage}\tabularnewline
\midrule
\endhead
\begin{minipage}[t]{0.38\columnwidth}\raggedright\strut
Salario mensual neto\strut
\end{minipage} & \begin{minipage}[t]{0.20\columnwidth}\raggedleft\strut
{1008,15}\strut
\end{minipage}\tabularnewline
\begin{minipage}[t]{0.38\columnwidth}\raggedright\strut
IRPF (9,64\%)\strut
\end{minipage} & \begin{minipage}[t]{0.20\columnwidth}\raggedleft\strut
115,68\strut
\end{minipage}\tabularnewline
\begin{minipage}[t]{0.38\columnwidth}\raggedright\strut
SS (30 \%)\strut
\end{minipage} & \begin{minipage}[t]{0.20\columnwidth}\raggedleft\strut
360,00\strut
\end{minipage}\tabularnewline
\begin{minipage}[t]{0.38\columnwidth}\raggedright\strut
Salario mensual bruto\strut
\end{minipage} & \begin{minipage}[t]{0.20\columnwidth}\raggedleft\strut
1200,00\strut
\end{minipage}\tabularnewline
\midrule
\begin{minipage}[t]{0.38\columnwidth}\raggedright\strut
\textbf{Total tres meses}\strut
\end{minipage} & \begin{minipage}[t]{0.20\columnwidth}\raggedleft\strut
3600\strut
\end{minipage}\tabularnewline
\bottomrule
\caption{Costes de personal.}
\end{longtable}

\textbf{Costes de \emph{hardware}:}

En este apartado se revisan los costes del \emph{hardware} desarrollado en el proyecto.

\begin{longtable}[]{@{}lr@{}}
\toprule
\begin{minipage}[b]{0.38\columnwidth}\raggedright\strut
\textbf{Concepto}\strut
\end{minipage} & \begin{minipage}[b]{0.20\columnwidth}\raggedleft\strut
\textbf{Coste \euro{}}\strut
\end{minipage}\tabularnewline
\midrule
\endhead
\begin{minipage}[t]{0.38\columnwidth}\raggedright\strut
PCB (10 prototipos)\strut
\end{minipage} & \begin{minipage}[t]{0.20\columnwidth}\raggedleft\strut
1,75\strut
\end{minipage}\tabularnewline
\begin{minipage}[t]{0.38\columnwidth}\raggedright\strut
Transporte\strut
\end{minipage} & \begin{minipage}[t]{0.20\columnwidth}\raggedleft\strut
7,11\strut
\end{minipage}\tabularnewline
\begin{minipage}[t]{0.38\columnwidth}\raggedright\strut
Componentes (para 2 uds)\strut
\end{minipage} & \begin{minipage}[t]{0.20\columnwidth}\raggedleft\strut
51,90\strut
\end{minipage}\tabularnewline
\begin{minipage}[t]{0.38\columnwidth}\raggedright\strut
Aduanas\strut
\end{minipage} & \begin{minipage}[t]{0.20\columnwidth}\raggedleft\strut
15,00\strut
\end{minipage}\tabularnewline
\begin{minipage}[t]{0.38\columnwidth}\raggedright\strut
Soldador\strut
\end{minipage} & \begin{minipage}[t]{0.20\columnwidth}\raggedleft\strut
15,00\strut
\end{minipage}\tabularnewline
\begin{minipage}[t]{0.38\columnwidth}\raggedright\strut
Estaño\strut
\end{minipage} & \begin{minipage}[t]{0.20\columnwidth}\raggedleft\strut
5,00\strut
\end{minipage}\tabularnewline
\begin{minipage}[t]{0.38\columnwidth}\raggedright\strut
Flux\strut
\end{minipage} & \begin{minipage}[t]{0.20\columnwidth}\raggedleft\strut
3,00\strut
\end{minipage}\tabularnewline
\begin{minipage}[t]{0.38\columnwidth}\raggedright\strut
Ordenador\strut
\end{minipage} & \begin{minipage}[t]{0.20\columnwidth}\raggedleft\strut
855,47\strut
\end{minipage}\tabularnewline
\midrule
\begin{minipage}[t]{0.38\columnwidth}\raggedright\strut
\textbf{Total}\strut
\end{minipage} & \begin{minipage}[t]{0.20\columnwidth}\raggedleft\strut
954,23\strut
\end{minipage}\tabularnewline
\bottomrule
\caption{Costes de personal.}
\end{longtable}



\textbf{Costes totales:}

\begin{longtable}[]{@{}lr@{}}
\toprule
\begin{minipage}[b]{0.38\columnwidth}\raggedright\strut
\textbf{Concepto}\strut
\end{minipage} & \begin{minipage}[b]{0.20\columnwidth}\raggedleft\strut
\textbf{Coste \euro{}}\strut
\end{minipage}\tabularnewline
\midrule
\endhead
\begin{minipage}[t]{0.38\columnwidth}\raggedright\strut
Personal\strut
\end{minipage} & \begin{minipage}[t]{0.20\columnwidth}\raggedleft\strut
3600 \euro{}\strut
\end{minipage}\tabularnewline
\begin{minipage}[t]{0.38\columnwidth}\raggedright\strut
\emph{Hardware}\strut
\end{minipage} & \begin{minipage}[t]{0.20\columnwidth}\raggedleft\strut
98,76 \euro{}\strut
\end{minipage}\tabularnewline
\midrule
\begin{minipage}[t]{0.38\columnwidth}\raggedright\strut
\textbf{Total}\strut
\end{minipage} & \begin{minipage}[t]{0.20\columnwidth}\raggedleft\strut
3698,76 \euro{}\strut
\end{minipage}\tabularnewline
\bottomrule
\caption{Costes totales.}
\end{longtable}

\textit{Nota: No se ha tenido en cuenta en el cálculo anterior el coste de los tutores.}

Teniendo en cuenta los datos anteriores, para hacer rentable el proyecto con un periodo de tiempo de un año, sería necesario vender tanto el \emph{hardware} como licencias de la aplicación. 

\subsection{Viabilidad legal}

A continuación se exponen todos los temas relacionados con las licencias del software y \emph{hardware} de proyecto, así como de la documentación e imágenes.

Tabla de licencias de las dependencias utilizadas:

\begin{longtable}[]{@{}llll@{}}
\toprule
\begin{minipage}[b]{0.18\columnwidth}\raggedright\strut
Dependencia\strut
\end{minipage} & \begin{minipage}[b]{0.10\columnwidth}\raggedright\strut
Versión\strut
\end{minipage} & \begin{minipage}[b]{0.49\columnwidth}\raggedright\strut
Descripción\strut
\end{minipage} & \begin{minipage}[b]{0.11\columnwidth}\raggedright\strut
Licencia\strut
\end{minipage}\tabularnewline
\midrule
\endhead
\begin{minipage}[t]{0.18\columnwidth}\raggedright\strut
Java\strut
\end{minipage} & \begin{minipage}[t]{0.08\columnwidth}\raggedright\strut
8.0.191\strut
\end{minipage} & \begin{minipage}[t]{0.49\columnwidth}\raggedright\strut
JDK\strut
\end{minipage} & \begin{minipage}[t]{0.11\columnwidth}\raggedright\strut
GPL\strut
\end{minipage}\tabularnewline
\begin{minipage}[t]{0.18\columnwidth}\raggedright\strut
USBtinLib\strut
\end{minipage} & \begin{minipage}[t]{0.08\columnwidth}\raggedright\strut
1.2.0\strut
\end{minipage} & \begin{minipage}[t]{0.49\columnwidth}\raggedright\strut
Librería para la conexión con el \emph{hardware}.\strut
\end{minipage} & \begin{minipage}[t]{0.11\columnwidth}\raggedright\strut
GPL\strut
\end{minipage}\tabularnewline
\begin{minipage}[t]{0.18\columnwidth}\raggedright\strut
JSSC\strut
\end{minipage} & \begin{minipage}[t]{0.08\columnwidth}\raggedright\strut
2.8.0\strut
\end{minipage} & \begin{minipage}[t]{0.49\columnwidth}\raggedright\strut
Librería para la búsqueda de puertos serie.\strut
\end{minipage} & \begin{minipage}[t]{0.11\columnwidth}\raggedright\strut
LGPL\strut
\end{minipage}\tabularnewline
\begin{minipage}[t]{0.18\columnwidth}\raggedright\strut
sqlite-jdbc\strut
\end{minipage} & \begin{minipage}[t]{0.08\columnwidth}\raggedright\strut
3.23.1\strut
\end{minipage} & \begin{minipage}[t]{0.49\columnwidth}\raggedright\strut
Librería para realizar la conexión con la base de datos.\strut
\end{minipage} & \begin{minipage}[t]{0.11\columnwidth}\raggedright\strut
Apache 2.0\strut
\end{minipage}\tabularnewline
\begin{minipage}[t]{0.18\columnwidth}\raggedright\strut
json\strut
\end{minipage} & \begin{minipage}[t]{0.08\columnwidth}\raggedright\strut
1.6\strut
\end{minipage} & \begin{minipage}[t]{0.49\columnwidth}\raggedright\strut
Librería para la importación y exportación de ficheros JSON.\strut
\end{minipage} & \begin{minipage}[t]{0.11\columnwidth}\raggedright\strut
JSON\strut
\end{minipage}\tabularnewline
\bottomrule
\caption{Dependencias del proyecto y sus licencias}
\end{longtable}


Teniendo en cuenta todas las licencias de las dependencias utilizadas, la licencia que más se ajusta a nuestro proyecto es la \emph{GNU General Public Licence v3.0} más conocida como (\emph{{GNU GPL}}.

Con el uso de esta licencia, nos exponemos a las siguientes condiciones:


\begin{longtable}[]{@{}lll@{}}
\toprule
\begin{minipage}[b]{0.20\columnwidth}\raggedright\strut
Derechos\strut
\end{minipage} & \begin{minipage}[b]{0.33\columnwidth}\raggedright\strut
Condiciones\strut
\end{minipage} & \begin{minipage}[b]{0.33\columnwidth}\raggedright\strut
Limitaciones\strut
\end{minipage}\tabularnewline
\midrule
\endhead
\begin{minipage}[t]{0.20\columnwidth}\raggedright\strut
Uso comercial\strut
\end{minipage} & \begin{minipage}[t]{0.33\columnwidth}\raggedright\strut
liberar código fuente.\strut
\end{minipage} & \begin{minipage}[t]{0.33\columnwidth}\raggedright\strut
Limitación de responsabilidad.\strut
\end{minipage}\tabularnewline
\begin{minipage}[t]{0.20\columnwidth}\raggedright\strut
Modificación\strut
\end{minipage} & \begin{minipage}[t]{0.33\columnwidth}\raggedright\strut
Nota sobre la licencia y copyright.\strut
\end{minipage} & \begin{minipage}[t]{0.33\columnwidth}\raggedright\strut
Sin garantías.\strut
\end{minipage}\tabularnewline
\begin{minipage}[t]{0.20\columnwidth}\raggedright\strut
Distribución\strut
\end{minipage} & \begin{minipage}[t]{0.33\columnwidth}\raggedright\strut
Modificaciones bajo la misma licencia.\strut
\end{minipage} & \begin{minipage}[t]{0.33\columnwidth}\raggedright\strut
\strut
\end{minipage}\tabularnewline
\bottomrule
\caption{Derechos y condiciones de la licencia GNU GPL.}
\end{longtable}





\apendice{Especificación de Requisitos}

\section{Introducción}

\section{Objetivos generales}

\section{Catalogo de requisitos}

\section{Especificación de requisitos}



\apendice{Especificación de diseño}

\section{Introducción}

En esta sección se expone cómo se han resuelto las especificaciones indicadas anteriormente.

Se definen los datos que va a utilizar la aplicación, el diseño de sus interfaces y su arquitectura.

\section{Diseño de datos}

\subsection{Base de datos}

La aplicación utiliza una base de datos SQLite para almacenar los proyectos. La estructura de la base de datos se compone de dos tablas:

\begin{itemize}
\item
\textbf{projects}: En esta tabla se almacenan los nombres de los proyectos. Dispone de un único campo llamado \emph{name}, el cual es clave primaria de la tabla.
\item
\textbf{data}: Esta tabla almacena los datos identificados en cada proyecto. Dispone de 4 campos:
\begin{itemize}
\item
\textbf{name}: Almacena la etiqueta con la que se ha identificado al dato.
\item
\textbf{ID}: Almacena la ID del \emph{frame} que contiene el dato identificado.
\item
\textbf{byte}: Almacena el \emph{byte} del frame que contiene el dato identificado.
\item
\textbf{projectName}: Clave foránea a la tabla \emph{projects}. Nos indica el proyecto al que pertenece el dato identificado.
\end{itemize}
\end{itemize}

A continuación se expone un diagrama relacional de la base de datos:

\imagen{databaseDiagram}{Diagrama relacional de la base de datos}
\newpage
\section{Diseño procedimental}

\imagen{processDiagram}{Diagrama de secuencia de la aplicación.}

Como podemos observar en el diagrama de secuencia, desde el puto de inicio del programa, podemos seleccionar tres ramas:
\begin{itemize}
\item
Comenzar un análisis
\item
Comenzar una monitorización.
\item
Comenzar una lectura de los datos en formato \emph{bruto}.
\end{itemize}

Para llegar a todas las opciones anteriores, es necesario haber configurado la interfaz con la que nos vamos a conectar al bus CAN.

Tanto la primera como la segunda línea del procedimiento hacen uso de la base de datos y de la gestión de proyectos, sin embargo la última rama no hace uso de la base de datos, ya que simplemente se encarga de monitorizar todos los datos en tiempo real en una tabla, sin realizar un almacenaje en disco.

\section{Diseño arquitectónico}

\subsection{Model-View-Controller (MVC)}

Aplicando este patrón de diseño, se utilizan 3 componentes, las vistas, los modelos y los controladores, de tal forma que se separa la lógica de la vista de la aplicación. Con el uso de este método, si realizamos una modificación en una parte de nuestro código, la otra parte no se ve afectada.

Por ejemplo, si modificamos la base de datos, solo modificaríamos el modelo encargado de esa acción, sin tener que modificar por ejemplo, la parte de la vista.

Está compuesto por tres componentes:

\textbf{Model:} Normalmente esta parte se encarga de los datos (no siempre). Esto puede ser por ejemplo, consultando a una base de datos para obtener información.

\textbf{View:} Componen la representación visual de los datos, es decir, todo lo que tenga que ver con la interfaz gráfica.

\textbf{Controller:} Es un mediador entre los modelos y las vistas. Se encarga de recibir las órdenes del usuario a través de la vista, realizar la petición de datos al modelo y devolverlo de nuevo a la vista para que sea mostrado al usuario.

\imagen{modeloMVC}{Modelo MVC.}
\newpage
\subsection{Arquitectura general}

A continuación se describe de manera general la arquitectura de la aplicación. Como se ha indicado anteriormente, se utiliza el modelo MVC para el funcionamiento interno de la aplicación.

Además de eso, cabe destacar la interacción con la base de datos, tanto de lectura y escritura, como la interacción con la interfaz CAN encargada de leer los datos que fluyen por el bus.

\imagen{arquitectura}{Arquitectura general del proyecto.}

\subsection{Diseño de paquetes}

Para conseguir una organización interna del proyecto, se agruparon las distintas funcionalidades de la aplicación en paquetes. De esta manera se consigue una aplicación modular y con funcionalidades independientes.

\imagen{paquetes}{Paquetes de la aplicación}

A continuación se describen cada uno de los paquetes:

\begin{itemize}
\item
\textbf{can}: En este paquete se almacena todo lo relacionado con la interacción con las interfaces CAN.
\item
\textbf{db}: En este paquete se almacena todo lo relacionado con la conexión a la base de datos, la obtención y la inserción de datos en la misma.
\item
\textbf{gui}: En este paquete se almacena todo lo relacionado con la interfaz gráfica. 
\item
\textbf{main}: En este paquete se encuentra el \emph{main} de la aplicación, así como el \emph{preloader}.
\item
\textbf{projects}: En ete paquete se almacena todo lo relacionado con la gestión de proyectos de la aplicación.
\item
\textbf{staticdata}: En este paquete se almacenan una serie de estructuras las cuales son utilizadas por distintos componentes del proyecto. En el se almacenan todos los datos que son representados por las gráficas de la aplicación.
\end{itemize}

\imagen{paquetesClases}{Diagrama de clases y paquetes de la aplicación.}

\imagen{paquetesClasesGui}{Diagrama de clases y paquetes de la interfaz gráfica.}



\subsection{Diseño de clases}

En este apartado se describe cada una de las clases que forman la aplicación y su funcionamiento.

\begin{landscape}
\imagenflotante{CRET}{Diagrama UML de la aplicación.}
\end{landscape}

\newpage
\section{Diseño de interfaces}

El diseño de las interfaces se ha llevado a cabo a través de la utilidad SceneBuilder. Se ha intentado seguir una estructura limpia y clara, para facilitar su interacción con el usuario.

A continuación se muestran algunas de las vistas de la aplicación:

\imagen{preloader}{Preloader de la aplicación.}

\imagen{analysis1}{Vista de análisis de la aplicación.}

\imagen{realTimeTag}{Vista de la monitorización de la aplicación.}

El siguiente diagrama nos muestra el flujo de navegabilidad por la aplicación.


\imagen{navegacion}{Diagrama de navegabilidad}.



\chapter[Documentación de programación]{Documentación técnica de programación}
\section{Introducción}

En este anexo se describe la documentación técnica de programación, incluyendo la preparación del entorno de desarrollo, la estructura de la aplicación y las configuraciones necesarias.

\section{Estructura de directorios}

El proyecto tiene la siguiente estructura:

\begin{itemize}
\item
\textit{/documentacion/:} Documentación del proyecto.
\item
\textit{/documentacion/img/:} Imágenes utilizadas en la documentación.
\item
\textit{/documentacion/javadoc/:} Documentación \emph{Javadoc}.
\item
\textit{/hardware/:} Esquemas del \emph{hardware} realizado en formato para KiCAD.
\item
\textit{/hardware/gerber/:} Ficheros \emph{gerber} para la producción del \emph{hardware}.
\item
\textit{/src/:} Código fuente de la aplicación Java.
\item
\textit{/resources/:} Recursos de la aplicación.

\end{itemize}


\section{Manual del programador}

El siguiente manual tiene como objetivo servir de referencia para futuros programadores. Se explica como montar el entorno de desarrollo, las dependencias necesarias así como su compilación y ejecución.

\subsection{Entorno de desarrollo}

Para el proyecto se necesita tener instaladas las siguientes dependencias:

\begin{itemize}
\item
Java JDK 8
\item
JavaFX
\item
Medusa\cite{javafx:medusa}
\item
Git
\item
SceneBuilder
\item
Eclipse
\end{itemize}

\subsection{Java JDK 8}

Es uno de los lenguajes más utilizados hoy en día. El uso del lenguaje Java nos permite la compatibilidad de la aplicación en distintas plataformas.

Se puede obtener del siguiente enlace: \url{https://www.oracle.com/technetwork/java/javase/downloads/jdk8-downloads-2133151.html}

\subsection{JavaFX}

Se trata de un conjunto de productos y tecnologías para la creación de interfaces gráficas en Java. Para la ejecución de estas aplicaciones solamente es necesaria la instalación de JRE en la máquina.

Se puede obtener del siguiente enlace:
\url{https://www.oracle.com/technetwork/java/javafx/overview/index.html}

\subsection{Git}

Para utilizar el repositorio en el que se encuentra la aplicación, es necesario el uso de Git. Este programa nos permitirá clonar el repositorio a nuestra máquina y empezar a trabajar con ello.

Podemos instalarlo en Windows habilitando el WSL (Windows Subsystem for Linux), y ejecutando el comando \emph{apt-get install git} con permisos de \emph{root}.

\imagen{gitClone}{Comando Git sobre WSL.}

\subsection{SceneBuilder}

Se trata de una interfaz gráfica creada para el desarrollo de interfaces gráficas en JavaFX. Sigue el modelo drag and drop, y toda la información es almacenada en un fichero FXML, un formato especial que extiende de XML.

\imagen{sceneBuilder}{SceneBuilder 10.0}

\subsection{Eclipse}

Para el desarrollo del código Java, se ha utilizado la IDE Eclipse. En ella se han importado las librerías correspondientes, e instalado plugins para el soporte de JavaFX. De esta manera, es posible llamar a SceneBuilder desde Eclipse.


\subsection{Obtención del código fuente}

El código fuente de la aplicación está hospedado en un repositorio Git en GitHub. Para obtener una copia de este código, se pueden seguir los siguientes pasos:

Abrir el WSL en Windows.

Posicionarse en el directorio en el que se desea copiar el código. (Si deseamos acceder al sistema de ficheros de Windows desde WSL, la ruta será /mnt/c).

Introducir el siguiente comando:

\textit{git clone https://github.com/amb0070/CRET.git}

Una vez finalizada la descarga del repositorio, se podrá acceder al código fuente y a los recursos del proyecto.

\subsection{Importar proyecto a Eclipse}

Una vez descargado el proyecto, abrimos Eclipse.

Hacemos \emph{click} en \emph{File} > \emph{Import...}

Seleccionamos \emph{Projects from File System or Archive}, y seleccionamos el directorio donde hemos clonado el repositorio.

De esta manera, queda importado el proyecto a nuestro \emph{workspace} de Eclipse.


\subsection{Importar librerías necesarias}

Para importar las librerías necesarias y utilizadas por el proyecto, podemos seguir los siguientes pasos:

Sobre la carpeta del proyecto, en el árbol izquierdo, hacer click derecho.

Seleccionar sub-menú \emph{Build Path} > \emph{Configure Build Path}.

Hacemos \emph{click} en el botón \emph{Add External JARs...} y seleccionamos los ficheros .jar de las librerías.


\section{Compilación, instalación y ejecución del proyecto}


\subsection{Compilación}
La compilación del proyecto se realiza desde Eclipse. 

Para exportarlo como un fichero .jar, podemos hacer click derecho sobre el árbol izquierdo del proyecto, y \emph{click} en \emph{Export...}.

Seleccionamos \emph{Java} > \emph{Runnable JAR file}.

Incluimos las librerías en el fichero .jar, y exportamos.

También es posible la compilación del proyecto con el uso de \emph{Maven}.
Para ello se ha creado un fichero \textit{pom.xml} en el cual se definen las dependencias con las cuales tiene que ser compilado el código y todas ellas son incluidas en el \textit{.jar} generado. De esta manera se asegura que el programa no tenga dependencias sin cumplir durante la ejecución en cualquier equipo.

\subsection{Ejecución}

No es necesario realizar una instalación de la aplicación. Simplemente tener instalado en la máquina el JRE con una versión igual o superior a la 1.8.0.

La ejecución de la aplicación no requiere de ningún elemento externo, pero si que es necesario tener el \emph{hardware} conectado a la máquina para su funcionamiento.

De igual manera, el proyecto puede ser ejecutado desde Eclipse, pero para su funcionamiento es necesario el \emph{hardware} que se ha diseñado.

\section{Producción del \emph{hardware}}

A continuación se expone el procedimiento para la producción del \emph{hardware} desarrollado.

El proyecto puede ser abierto con KiCAD de forma que tenemos acceso a la modificación de los esquemas y del diseño de la PCB.

El flujo básico de trabajo en KiCAD, a rasgos muy generales sería el siguiente:
\begin{itemize}
\item
\textbf{Esquema}: Desarrollo del esquema del circuito. En éste esquema indicaríamos los componentes a utilizar y su conexión entre ellos.
\imagen{esquemaAnexo}{Ejemplo del diseño de uno de los elementos en el esquema.}

\item
\textbf{Asignación de huellas}: El siguiente paso en el desarrollo sería asignar a cada uno de los elementos del esquema el \emph{footprint}, es decir, la forma y tamaño del componente que se va a incluir en la \emph{PCB}.

Para ello, podemos hacer \emph{click} en el siguiente botón:

\imagen{iconoFootprint}{Icono para asignar las huellas al esquema.}

En la siguiente ventana tenemos una relación entre el nombre del componente y el \emph{footprint} asignado.

\imagen{asignamiento}{Relación entre los componentes y su \emph{footprint}.}

Analizando el primer elemento marcado en la lista, podemos obtener la siguiente información:

\begin{itemize}
\item
\textit{C1}: Hace referencia al nombre del elemento.
\item
\textit{100nF}: Hace referencia a la capacidad del elemento. En este caso al ser un condensador, la medida está indicada en nano faradios.
\item
\textit{Capacitor SMD:} Hace referencia al tipo de componente. En este caso un condensador para montaje sobre superficie (\emph{SMD - Surface Mount Device}).
\item
\textit{1206}: Este dato hace referencia a la tecnología de montaje superficial. Se podría decir que nos indican el tamaño del encapsulado del componente.
\item
\textit{Pad1.42x.75mm HandSolder}: En este caso, nos indica el tamaño de los \emph{pads}, es decir, de la superficie en la que es soldado el componente. En este caso tienen un tamaño un poco más grande de lo normal para facilitar la soldadura a mano (\emph{HandSolder}).
\end{itemize}

\item
\textbf{Netlist:} Una vez concluido el diseño del esquema y su asignación de huellas, deberíamos de generar el fichero \emph{netlist}. Este fichero será el que posteriormente se importe al diseñador de \emph{PCBs} y el cual contiene todos los elementos del esquema y sus conexiones.

En \emph{KiCAD} podemos generarlo con el siguiente botón:

\imagen{iconoNetlist}{Icono para generar el fichero \emph{netlist}.}

Importamos dicho fichero en el editor de \emph{PCBs} y obtendremos algo parecido a esto:

\imagen{firstFootprint}{Diseño inicial de la \emph{PCB}}

Insertando las vías y \emph{enrutando} todo el diseño, obtenemos el siguiente resultado:

\imagen{finalFootprint}{Diseño final de la \emph{PCB}.}

\end{itemize}

Con el diseño concluido, generamos los ficheros \emph{Gerber} para su producción.

En este caso, los prototipos han sido producidos en:

\url{https://jlcpcb.com/}.

El listado de componentes necesarios para la producción del \emph{hardware} son los siguientes:


\begin{center}
 \begin{tabular}{||c c c||} 
 \hline
 Cantidad & Componente & Elemento \\ [0.8ex] 
 \hline\hline
 2 & PIC18F14K50T-I/SO & IC1-1, IC1-2 \\ 
 \hline
 2 & MCP2515T-I/SOCT-ND & IC2-1, IC2-2\\
  \hline
 1 & GL850G & U2\\
   \hline
 1 & SY6280 & U1\\
 \hline
 2 & MCP2551-I/SN-ND & IC3-1, IC3-2\\
 \hline
 2 & CRYSTAL 24.0000MHZ 18PF SMD & Q1, Q2\\
 \hline
 2 & FERRITE BEAD 600 OHM 1206 1LN & L1, L2, L3, L4\\
  \hline
 2 & LED RED DIFFUSED 1206 SMD & LED1, LED2\\
  \hline
 1 & CAP CER 1UF 10V X7R 1206 & C17\\
   \hline
 2 & CAP CER 33PF 50V C0G/NP0 1206 & C13, C14\\
  \hline
 1 & CONN RCPT USB2.0 MICRO B SMD R/A & J7\\
  \hline
 2 & CONN D-SUB PLUG 9POS R/A SOLDER & CAN1, CAN2\\
  \hline
 4 & CAP CER 22PF 10V C0G/NP0 1206 & C10, C11, C18, C19\\
  \hline
 3 & CAP CER 10UF 10V X7R 1206 & C7, C15, C22\\
  \hline
 2 & CAP CER 0.1UF 10V X7R 1206 & C8, C9\\
  \hline
 15 & CAP CER 100UF 6.3V X5R 1206 & C(1-6,12,16,20,21,23-25),CP(3,4)\\
  \hline
 2 & CAP CER 4.7UF 10V X5R 1206 & CP1, CP2\\
   \hline
 2 & RES 47K OHM 5\% 1/4W 1206  & R9, R12\\
   \hline
 1 & RES 5.1K OHM 5\% 1/4W 1206  & R8\\
   \hline
 3 & RES 10K OHM 0.1\% 1/8W 1206  & R7, R10, R11\\
   \hline
 1 &  CRGCQ 1206 680R 1\% & R6\\
   \hline
 2 &  RES SMD 120 OHM 5\% 1/2W 1206  & R4, R15\\
    \hline
 8 &  RES 1K OHM 1\% 1/2W 1206  & R2,R3,R5,R13,R14,R16-R18\\
 \hline
 1 &  RES SMD 4.7K OHM 1\% 1/4W 1206 & R1 \\ 
  \hline
 1 &  DIODE 1206 & D1 \\[1ex] 
 \hline
\end{tabular}
\end{center}


Todos estos componentes han sido comprados en \url{https://www.digikey.es/} aunque pueden ser obtenidos en otras tiendas.

El formato de todas las resistencias, condensadores y ferritas ha sido SMD 1206, debido a su reducido tamaño pero a la vez suficiente para soldarlo a mano.

Una vez disponemos de la PCB y de los componentes, es necesario realizar un paso previo antes de soldarlos. Es necesario programar los PIC con el \emph{bootloader} para posteriormente cargar el \emph{firmware} en los mismos.

Para ello, teniendo en cuenta el \emph{datasheet}\cite{datasheet} del microcontrolador, podemos realizar la siguiente conexión:


\imagen{conexionPickit}{\emph{Pinout} del \emph{PICKit3}\cite{PICKIT3:pinout}.}


\imagen{cargaFirmware}{Cargando \emph{firmware} en el PIC.}


Una vez detectado el dispositivo por el equipo, es necesario conectarse a el, e indicar el fichero .hex\cite{hexfile} que queremos cargar, en este caso, el \emph{bootloader} referenciado en la página oficial.

Es importante configurar el programador para que alimente el microcontrolador, sino la carga del \emph{bootloader} no será posible. Para ello, hay que acceder al modo avanzado del programador.

\imagen{pickit3}{Configuración avanzada de \emph{MPLAB}\cite{mplab}.}

Seleccionamos la opción para que el programador alimente el controlador.

Ya podemos grabar el \emph{bootloader} en el microcontrolador.

Una vez realizados estos pasos, ya se pueden soldar todos los componentes en la placa.

Para hacerla funcionar es necesario un paso más. Cargar el \emph{firmware}.

Para ello, conectamos la placa al USB del equipo, con el \emph{jumper} JP3 o JP4 (cada uno es para su microcontrolador), y ejecutamos la siguiente utilidad:

mphidflash - \url{https://github.com/ApertureLabsLtd/mphidflash}

\imagen{mphidflash}{Salida del comando \emph{mphidflash}\cite{mphidflash}.}


Una vez completada la carga, podemos comprobar el dispositivo y estaría listo para su uso.

\apendice{Documentación de usuario}

\section{Introducción}

\section{Requisitos de usuarios}

\section{Instalación}

\section{Manual del usuario}





\bibliographystyle{plain}
\bibliography{bibliografiaAnexos}

\end{document}
